\section{Limitaciones, criterio de ingeniería y direcciones futuras}

El summary y el future work del final de la clase son mesurados. La contribución es una perspectiva nueva, sencilla, estable y extensible. No anuncian que la language modeling haya reemplazado al RL. Esta sección convierte las limitaciones más importantes del Q\&A de la clase en puntos de verificación aplicables.

\subsection{Síntesis de la clase: por qué el método merece atención}

El valor de Decision Transformer tiene tres fuentes. Primera: usa conditional sequence modeling para unificar return, state y action. Segunda: entrena con un objetivo supervisado estable, sin depender explícitamente de una value function ni de policy gradient. Tercera: obtiene resultados competitivos en varios tipos de offline benchmark, y los experimentos de context, sparse reward y critic muestran el papel de la sequence history.

\lecturefigure{slide-22-summary.jpg}{Síntesis de la conferencia original: sencillo, estable, fácil de entrenar y con buen desempeño en varios tipos de configuraciones de RL.}{Video oficial de Stanford Online, 1:08:56.}

\begin{importantbox}{Principio de proporcionalidad de la evidencia}
Se puede afirmar que Decision Transformer demuestra que el sequence modeling es una offline-RL formulation viable y sólida. No se puede afirmar que demuestre que todo el RL deba reemplazarse por next-token prediction, ni que basta con agrandar el Transformer para resolver exploration, safety y distribution shift.
\end{importantbox}

\subsection{El Online RL no se obtiene con solo agregar ruido}

En el offline setting, los datos de entrenamiento ya existen. El online RL tiene que recopilar datos mientras mejora la policy. En el Q\&A de la clase se preguntó si se podía adoptar directamente epsilon-greedy o Boltzmann sampling. El ponente respondió que la evidencia preliminar indica que la exploration recipe del RL tradicional no se puede trasladar tal cual. Tanto el exploration/exploitation balance como la forma en que las trayectorias erróneas entran al context requieren un diseño nuevo.

\begin{warningbox}{Tres problemas adicionales al pasar al modo en línea}
\begin{enumerate}
  \item \textbf{Cold start}: si no hay trajectory de alto retorno, el target conditioning no tiene una conducta correspondiente;
  \item \textbf{Non-stationary dataset}: las actualizaciones de la policy cambian continuamente la distribución de los datos;
  \item \textbf{Sequence contamination}: el historial de baja calidad que generan las acciones de exploración afecta las predicciones posteriores y no puede tratarse simplemente como iid augmentation.
\end{enumerate}
\end{warningbox}

\teachervoice{Nota del apunte: el ponente no afirmó que epsilon-greedy esté destinado a fallar. Dijo que «devil lies in the detail»: las técnicas existentes de RL en línea no son directamente transferable. La pregunta de investigación correcta es cómo diseñar la recopilación de datos, el condicionamiento por objetivo y el sequence replay, no solo ajustar una sampling temperature.}

\subsection{Discount factor y context limitado}

Los experimentos originales usan principalmente el finite-horizon return con $\gamma=1$, pero Decision Transformer es compatible con el discounting. Basta con usar $\widehat{R}^{(\gamma)}_t$ de forma consistente en el preprocesamiento de datos y en el rollout. Sin embargo, el discounting resuelve el peso de las reward lejanas y no equivale a memoria. Aunque $\gamma$ sea razonable, un $K$ limitado puede impedir que el modelo vea los eventos tempranos necesarios para completar la tarea.

\begin{importantbox}{Discount y context responden a preguntas distintas}
$\gamma$ determina el peso numérico de las reward lejanas; $K$ determina cuánta información histórica puede consultar el modelo. Un $\gamma$ menor no recupera una observation clave que se truncó, y un $K$ mayor tampoco sustituye una definición estable del return de horizonte infinito.
\end{importantbox}

\subsection{Relación con CQL, pessimism y el hybrid objective}

CQL (Conservative Q-Learning) contrarresta la offline overestimation reduciendo el valor estimado de las action que están fuera de los datos. En el Q\&A de la clase se preguntó si se podía agregar pessimism o un Q-learning objective a Decision Transformer. La respuesta es que en principio sí. Sin embargo, la pregunta científica del trabajo original era precisamente hasta dónde se puede llegar con un sequence objective sencillo, sin esos componentes tradicionales.

\begin{knowledgebox}{Tres direcciones hybrid}
\begin{enumerate}
  \item \textbf{Sequence actor + pessimistic critic}: DT genera acciones candidatas y un value model conservador las filtra;
  \item \textbf{Return model + uncertainty}: se modelan conjuntamente la alcanzabilidad del desired return y la confianza en él;
  \item \textbf{Online sequence replay}: los datos de contexto se mantienen y se actualizan continuamente según trajectory quality, novelty y safety.
\end{enumerate}
Estas direcciones aumentan la complejidad del algoritmo y obligan a volver a verificar la estabilidad y la contribución causal.
\end{knowledgebox}

\subsection{Trabajo futuro: multimodal, multitask, multi-agent}

La future-work slide del ponente trata el modelado de secuencias como una interfaz unificada. La multimodality permite que el modelo procese al mismo tiempo experience en línea y fuera de línea, imágenes y texto. El multitask requiere task-conditioned goal más ricos o lenguaje. El multi-agent incorpora a la secuencia conjunta las acciones y la observation de los otros agent. Cada dirección amplía el token space y, con él, los problemas de credit, coordination y safety.

\lecturefigure{slide-23-future-work.jpg}{Direcciones futuras: multimodality, multitask y multi-agent sequential decision making.}{Video oficial de Stanford Online, 1:10:46.}

\begin{knowledgebox}{Lectura de la figura: una interfaz de secuencia unificada no implica una dificultad unificada}
Convertir todas las modalidades en tokens puede unificar la pila de software, pero no unifica por sí solo la reward semantics, la action latency, la agent identity ni las restricciones de seguridad. Un multimodal model necesita sincronización temporal; un multitask model necesita evitar la confusión de objetivos; un multi-agent model, además, debe manejar oponentes no estacionarios y la credit assignment conjunta.
\end{knowledgebox}

\subsection{Recursos de la conferencia original}

La slide original indica la página del proyecto, el artículo y el código, y cada uno cumple una función distinta. La página del proyecto ofrece visualizaciones y una explicación breve. El artículo define el método y los experimentos. El código muestra los detalles del preprocesamiento de datos, el sequence sampling, la return scale y la evaluación. Para reproducir los resultados hay que verificar contra el código y la configuración original de los datos, no reescribir a partir del diagrama de arquitectura.

\lecturefigure{slide-24-useful-links.jpg}{Página del proyecto, artículo y código de Decision Transformer citados en la conferencia original.}{Video oficial de Stanford Online, 1:11:30.}

\begin{importantbox}{Lista de verificación para la reproducción}
\begin{enumerate}
  \item Verificar la reward normalization, la return scale y la elección del target-return;
  \item registrar la dataset collection policy, la trajectory length y la proporción de éxitos;
  \item alinear la context length, el sequence sampling y la padding mask;
  \item reportar a la vez la loss de entrenamiento, el realized return, el target sweep y la varianza entre varias semillas aleatorias;
  \item comparar al menos BC, percent BC, un value-based offline baseline y una context ablation.
\end{enumerate}
\end{importantbox}

\subsection{Resumen de la sección}

Los límites de Decision Transformer importan tanto como sus contribuciones. Ofrece un offline sequence objective estable y unificado. Aun así, la online exploration, las conductas que exceden el soporte de los datos, el context limitado, el discounting y la conservative policy improvement todavía deben resolverse por separado.
